\chapter{Ingrown Hair (Pseudofolliculitis Barbae)}

\section*{Definition}
An ingrown hair is a hair that curls back or grows sideways into the skin, causing inflammation. It most commonly affects the beard area, neck, legs, and bikini area. It is particularly common in people with curly or coarse hair.

\section*{What Happens in Your Body}
When a hair is shaved, the sharp tip can re-enter the skin as it grows. Curly hair is more prone to curling back into the follicle. The body mounts an inflammatory response, causing redness, papules, pustules, and sometimes hyperpigmentation or scarring. It is a foreign-body inflammatory reaction.

\section*{Causes}
\begin{itemize}
\item Shaving (most common cause)
\item Tweezing or waxing
\item Tight clothing (bikini area)
\item Curly or coarse hair type
\item Close shaving (multi-blade razors)
\item Shaving against the direction of hair growth
\item Stretching the skin while shaving
\item Dry shaving or inadequate lubrication
\end{itemize}

\section*{When to See a Doctor}
\begin{itemize}
\item Small, raised, red bumps (papules) in areas of hair removal
\item Pustules that may look like acne
\item Pain or tenderness
\item Itching
\item Hyperpigmentation in chronic cases
\item Curved hair visible under the skin
\item Razor bumps (pseudofolliculitis barbae) on the neck and beard area
\end{itemize}

\section*{Related Symptoms}
\begin{itemize}
\item Folliculitis (infection of hair follicles)
\item Acne vulgaris
\item Razor burn
\item Pseudofolliculitis (razor bumps)
\item Hyperpigmentation (post-inflammatory)
\end{itemize}

\section*{Self-Care / Home Management}
\begin{itemize}
\item Stop shaving the affected area until healed
\item Use a warm compress to bring the hair to the surface
\item Gently extract the embedded hair with a sterile needle
\item Shave with the grain, not against it
\item Use a single-blade razor
\item Apply warm water and shaving cream before shaving
\item Exfoliate gently with a soft brush or washcloth
\item Use electric clippers instead of a razor
\end{itemize}

\section*{How Your Doctor Will Evaluate You}
\begin{itemize}
\item Clinical examination of affected skin
\item Rule out other causes of folliculitis (bacterial, fungal)
\end{itemize}

\section*{Treatment Options}
\begin{itemize}
\item Discontinue hair removal temporarily
\item Apply topical antibiotics if secondarily infected
\item Topical retinoids (tretinoin) to normalize follicular keratinization
\item Topical corticosteroids for inflammation
\item Laser hair removal for chronic, recurrent cases
\item Electrolysis as a permanent solution
\item Consider chemical depilatories instead of shaving
\end{itemize}

\section*{References}

\begin{thebibliography}{99}
\bibitem{perry-pfb} Perry PK, et al. Pseudofolliculitis barbae: a clinical review. \emph{J Am Acad Dermatol}. 2023;88(3):501-512.
\bibitem{uptodate-pfb} Alexis A, et al. Pseudofolliculitis barbae. In: UpToDate, Post TW (Ed), UpToDate, Waltham, MA; 2025.
\bibitem{kindred} Kindred C, et al. Pseudofolliculitis barbae: treatment with laser hair removal. \emph{Dermatol Surg}. 2022;48(5):542-549.
\bibitem{james-wd} James WD, et al. \emph{Andrews' Diseases of the Skin}. 13th ed. Elsevier; 2023. Chapter 29: Pseudofolliculitis.
\bibitem{bridgeman} Bridgeman D, et al. Medical management of pseudofolliculitis barbae. \emph{Am J Clin Dermatol}. 2024;25(2):209-223.
\bibitem{gray-gray} Gray J, et al. Razor bumps: pathophysiology and prevention. \emph{Int J Dermatol}. 2023;62(8):987-995.
\end{thebibliography}

\clearpage
